
% Para o valor absoluto. Copiado das memorias de alguén
\providecommand{\abs}[1]{\lvert#1\rvert}

% novo comando pa que sea bonito referenciar cousas
\newcommand{\referencia}[1]{[\textbf{\ref{#1}}]}

% :FACER: diferencia entre esto e \DeclasreMathOperator
% Comandos para operadored diferenciales usuales
% Véxase: https://tex.stackexchange.com/a/67529/326141
\newcommand{\gradiente}{\operatorname{grad}}
\newcommand{\gr}[1]{\gradiente\left(#1\right)}

\newcommand{\diverxencia}{\operatorname{div}}
\newcommand{\dv}[1]{\diverxencia\left(#1\right)}

\newcommand{\rotacional}{\operatorname{rot}}
\newcommand{\rt}[1]{\rotacional\left(#1\right)}

\newcommand{\orden}{\operatorname{Ord}}
\newcommand{\ord}[1]{\ensuremath{\orden\left(#1\right)}}

\newcommand{\subgrupo}{\ensuremath{\le}}
\newcommand{\subgrupopropio}{\ensuremath{<}}
\newcommand{\subgruponormal}{\ensuremath{\trianglelefteq}}
\newcommand{\subgruponormalpropio}{\ensuremath{\triangleleft}}
\newcommand{\combinacion}[1]{\ensuremath{\left<#1\right>}}

% Espazos de matrices
% :FACER: mais tipos, con argumentos opcionais de letras n,m
\newcommand{\Mnn}[1]{\mathcal{M}_{n\times n}(\mathbb{#1})}

% Grupos
% :FACER: mais tipos, mellores nomes, argumento opcional pa mostrar o corpo
% \newcommand{\GL}[1]{\operatorname{GL}\ensuremath{(#1)}}
\newcommand{\SL}[1]{\operatorname{SL}\ensuremath{(#1)}}
\newcommand{\SU}[1]{\operatorname{SU}\ensuremath{(#1)}}
\newcommand{\SO}[1]{\operatorname{SO}\ensuremath{(#1)}}
\newcommand{\U}[1]{\operatorname{U}\ensuremath{(#1)}}
\newcommand{\Ort}[1]{\operatorname{O}\ensuremath{(#1)}}

% Alxebras
\newcommand{\gl}[1]{\operatorname{\mathfrak{gl}}\ensuremath{(#1)}}

% :FACER: facer o resto con esto de ifthen
% https://stackoverflow.com/a/7314951
\RequirePackage{ifthen}
\newcommand{\GL}[2][]{
    \ifthenelse { \equal {#1} {}}
        {\def\temp{\operatorname{GL}\ensuremath{(#2)}}}
        {\def\temp{\operatorname{GL}\ensuremath{(#2,\mathbb{#1})}}}
    \temp
}


% para vectores
% \newcommand{\ve}[1]{\mathbf{#1}}

\endinput
